% Fig. 4 按类 ΔF1 对真实样本 10-NN 纯度（原划分，10 种子；点 = 4 个检测器的均值，竖线 = 4 个检测器的最小–最大）
% 数据：data/fig4_<来源>.dat（由 prep_figs.py 生成）；编译：pdflatex fig4_perclass.tex
\documentclass[10pt]{standalone}
% 全部论文图共用的样式（颜色固定跟随来源；字号按 elsarticle 正文 10pt 设定）
% 颜色来自经过色觉校验的分类调色板，顺序固定：
%   类别权重 = 1 蓝，ROS = 2 橙，SMOTE = 3 青，Qwen = 4 黄，Llama = 5 品红；不增强 = 中性灰
\usepackage{pgfplots}
\usepgfplotslibrary{groupplots}
\pgfplotsset{compat=1.18}
\usetikzlibrary{plotmarks,calc,patterns,arrows.meta,positioning,shapes.geometric,fit,backgrounds}
\definecolor{cCW}{HTML}{2A78D6}
\definecolor{cROS}{HTML}{EB6834}
\definecolor{cSMOTE}{HTML}{1BAF7A}
\definecolor{cQWEN}{HTML}{EDA100}
\definecolor{cLLAMA}{HTML}{E87BA4}
\definecolor{cNONE}{HTML}{52514E}
\definecolor{cGRID}{HTML}{E4E3DF}
\definecolor{cAXIS}{HTML}{8A8984}
\definecolor{cTEXT}{HTML}{2B2A28}
\pgfplotsset{
  paper/.style={
    font=\footnotesize, color=cTEXT,
    axis line style={cAXIS, line width=0.4pt},
    tick style={cAXIS, line width=0.4pt},
    major tick length=2pt,
    grid style={cGRID, line width=0.4pt},
    ymajorgrids,
    tick label style={font=\scriptsize},
    label style={font=\footnotesize},
    title style={font=\footnotesize\bfseries, yshift=-2pt},
    legend style={font=\scriptsize, draw=none, fill=none, /tikz/every even column/.append style={column sep=8pt}},
    legend cell align=left,
  },
  errbar/.style={error bars/.cd, y dir=both, y explicit,
    error bar style={line width=0.5pt, cTEXT!75}, error mark options={rotate=90, mark size=1.4pt, line width=0.5pt, cTEXT!75}},
}

\begin{document}
\begin{tikzpicture}
\pgfplotsset{
  discard if not/.style 2 args={x filter/.append code={%
      \edef\tempa{\thisrow{#1}}\edef\tempb{#2}%
      \ifx\tempa\tempb\else\def\pgfmathresult{inf}\fi}},
  f4/.style={paper, xmajorgrids,
    xmin=-0.03, xmax=1.03, xtick={0,0.5,1}, xticklabels={0, 0.5, 1},
    ymin=-0.33, ymax=0.23, ytick={-0.3,-0.2,...,0.2},
    scaled y ticks=false,
    yticklabel style={/pgf/number format/fixed, /pgf/number format/fixed zerofill, /pgf/number format/precision=1, /pgf/number format/showpos},
    extra y ticks={0}, extra y tick labels={},
    extra y tick style={grid=major, major grid style={cTEXT!70, line width=0.6pt}},
    error bars/y dir=both, error bars/y explicit,
    error bars/error bar style={line width=0.6pt},
    error bars/error mark=none},
  ds0/.style={only marks, mark=*, mark size=1.9pt, discard if not={dsid}{0}},
  ds1/.style={only marks, mark=square*, mark size=1.7pt, discard if not={dsid}{1}},
  ds2/.style={only marks, mark=triangle*, mark size=2.3pt, discard if not={dsid}{2}},
}
\newcommand{\fourpanel}[2]{% #1 = 来源键名，#2 = 颜色
  \addplot[ds0, #2, mark options={fill=#2, draw=white, line width=0.3pt}, error bars/error bar style={#2!75!black}] table[x=real_purity, y=mean, y error plus=eplus, y error minus=eminus]{data/fig4_#1.dat};
  \addplot[ds1, #2, mark options={fill=#2, draw=white, line width=0.3pt}, error bars/error bar style={#2!75!black}] table[x=real_purity, y=mean, y error plus=eplus, y error minus=eminus]{data/fig4_#1.dat};
  \addplot[ds2, #2, mark options={fill=#2, draw=white, line width=0.3pt}, error bars/error bar style={#2!75!black}] table[x=real_purity, y=mean, y error plus=eplus, y error minus=eminus]{data/fig4_#1.dat};
  \draw[cAXIS, dashed, line width=0.5pt] (axis cs:0.5,-0.33) -- (axis cs:0.5,0.23);
}
\begin{groupplot}[
  group style={group size=5 by 1, horizontal sep=0.32cm, y descriptions at=edge left},
  width=3.95cm, height=5.0cm,
]
\nextgroupplot[f4, title={Class weight}, ylabel={$\Delta$F1 vs.\ no augmentation (per class)}, ylabel style={yshift=-3pt}]
  \fourpanel{cw}{cCW}
\nextgroupplot[f4, title={ROS}, xlabel={10-NN class purity of real training samples}, xlabel style={xshift=2.05cm}]
  \fourpanel{ros}{cROS}
\nextgroupplot[f4, title={SMOTE}]
  \fourpanel{smote}{cSMOTE}
  % CIC Web-SQLi 的 SMOTE 均值 −0.445 超出纵轴范围：在底部用箭头标出
  \draw[-{Latex[length=3pt]}, cSMOTE!75!black, line width=0.6pt] (axis cs:0.073,-0.24) -- (axis cs:0.073,-0.322);
  \node[font=\tiny, anchor=west, align=left, cTEXT, inner sep=1pt] at (axis cs:0.10,-0.285) {$-0.45$\\CIC Web-SQLi};
\nextgroupplot[f4, title={LLM (Qwen)}]
  \fourpanel{qwen}{cQWEN}
\nextgroupplot[f4, title={LLM (Llama)}]
  \fourpanel{llama}{cLLAMA}
\end{groupplot}
% 数据集图例（形状）
\node[anchor=south, font=\scriptsize] at ($(group c3r1.north)+(0,0.78cm)$) {%
  \tikz[baseline=-0.6ex]\fill[cNONE] (0,0) circle (1.9pt);~CSE-CIC-IDS2018\qquad
  \tikz[baseline=-0.6ex]\fill[cNONE] (-1.7pt,-1.7pt) rectangle (1.7pt,1.7pt);~NF-UNSW-NB15-v3\qquad
  \tikz[baseline=-0.6ex]\fill[cNONE] (-2pt,-1.6pt) -- (2pt,-1.6pt) -- (0,2pt) -- cycle;~CIC-IDS2017};
\end{tikzpicture}
\end{document}
